\chapter{Literature Review}\label{chap:literature}

Policy discussions regarding growth of developing countries often revolve around the role of informal economy. Informality remains a key concern of the Decent Work Agenda and the 2030 Agenda for Sustainable Development. The informality rate is an integral part of the indicators selected to measure progress towards the Sustainable Development Goals, as SDG indicator 8.3.1. It was speculated that with economic growth informal economies would diminish giving way to formalization. On the other hand formal-informal linkages have deepened. This dissertation contributes to understanding this puzzle by examining whether such linkages reflect efficient technological complementarities (the "displacement" view) or extractive power asymmetries (the "adverse incorporation" view).

This literature review is organized around four interconnected themes. First, the dissertation traces the theoretical models from dualist view of informality to structuralist perspectives that emphasize the functional role of informal labor in global value chains. Second, a review empirical evidence on formal-informal linkages through subcontracting and outsourcing, with particular attention to the Indian manufacturing context is carried out. Third, the literature review explores studies on labor institutions and their enforcement, which forms a crucial the theoretical foundation for identification strategy later. Finally, I discuss methodological approaches to estimating bargaining power and surplus distribution in vertically-linked production systems.

\section{Theoretical Frameworks: From Dualism to Structural Heterogeneity}

\subsection{Dualism of Production}
\autocite{Lewis1954} differentiated the subsistence
sector from the capitalist sector by articulating it as that part of the economy which is not using reproducible capital. Compared to the capitalist sector in the subsistence sector output per head is lower because it is not fructified by capital.
Subsequent extensions made by Harris and Todaro (1970) and Rauch (1991) helped in the development of the classical view of informality. According to the dual model developing economies are characterized by a bifurcated structure: a modern, capital-intensive formal sector coexisting with a traditional, low-productivity informal sector. 

\autocite{Lewis1954} also posited that economic development proceeds through the gradual transfer of surplus labor from the traditional to the modern sector, with the process continuing until the labor surplus is exhausted and wages begin to rise across the economy.
This perspective treats informality as a transitory phenomenon which will be replaced by the modern sphere of economic activity with development. Subsequent refinements incorporated migration costs (Harris and Todaro, 1970), credit constraints (Evans and Jovanovic, 1989), and regulatory barriers (Rauch, 1991), but retained the core prediction: as institutional quality improves and capital deepens, informal firms either formalize or exit, unable to compete with the superior productivity of formal enterprises.

The displacement view finds empirical support in large productivity gaps between formal and informal firms. La Porta and Shleifer (2014) document that informal firms in developing countries exhibit productivity levels 30-70\% below their formal counterparts, lack access to credit and modern technology, and rarely transition to formality. They argue that informal firms survive by evading taxes and regulations, not through genuine competitive advantage.

\textit{Dual economy scholars correctly categorize
informal enterprises as primarily small-scale and undercapitalized (in terms of
informal enterprises as primarily small-scale and undercapitalized (in terms of
both human and financial terms); they have failed to create a viable explanation
for the persistence of informal economic activity except as a distortion caused
by inefficiencies in the institutional or political structure of countries. Explanations about persistence come from sociologists or policy scholars, and I con
sider their contributions next.}

 \subsection{The Structuralist View of Production}
\textit{K. Hart (1973), based on his work in Accra, Ghana, noted that Lewis’s
model didn’t hold. Instead of finding an endless stream of labor flowing from the rural, low-productivity sector into the urban and formal sector, he
found a symbiotic dual urban economy: a limited formal sector and a thriving
urban informal sector, as well as a third, criminal sector. Hart actually coined
the term informal economy, and he noted that the two economies had quite a
number of bidirectional flows and interactions as individuals moved between
sectors, but that the proximity of the low-productivity subsistence sector did
not fuel investment and productivity in the formal sector. Instead, he found
an informal sector that helped many buffer the ups and downs of employment
in the formal economy.}

\begin{quote}
	In contrast with sectoral and dualistic approaches we maintain the need to conceptualize informality as a particular relationship in the organization of production. In other words, a central idea of our study is that informality is a feature of labor-capital relations. \parencite[pp. 78--79]{Fortuna1989}
\end{quote}

A fundamentally different perspective emerges from structuralist political economy, which views informality not as a relic of underdevelopment but as a functional requirement of contemporary capitalism. Drawing on dependency theory and world-systems analysis, scholars such as Portes, Castells, and Benton (1989), Sassen (1996), and Meagher (2010) argue that informal labor is systematically incorporated into formal production networks on highly unequal terms—a process termed "adverse incorporation." Sanyal (2007) distinguishes between  "need economy" activities (subsistence-level production for survival) and "formal economy" production that is integrated into capitalist accumulation. Crucially, he demonstrates how capital strategically deploys informal labor to extract surplus value while externalizing social reproduction costs. Rather than competing with formal firms, informal producers serve as subordinated suppliers, bearing volatility and regulatory risk while formal firms capture profits.

This perspective finds resonance in the global value chain (GVC) context, where scholars have shown that post-globalization the structure of global capitalism is not a uniform march toward formality. \textcite{milbergwinkler} by analyzing the relationship between offshoring and labor share in the U.S. and other OECD countries conclude that offshoring reduces the labor share, especially in countries with weaker labor market institutions. In addition, for developing economies participation in GVCs does not automatically lead to higher wages or better labor standards even when the economy is upgrading, for example in the case of Mexico. 

The GVCs that have emerged therefore present the case of powerful lead firms outsourcing to small, dependent suppliers in developing countries. In order to meet the demanding requirements (low cost, fast turnaround, flexibility) while bearing high risks and thin margins, these suppliers functionally depend on informal practices \parencite{gereffi2005governance}. While Gereffi et al. (2005) provided the structural logic for why global capitalism fragments production, \textcite{barrientos2011economic} reveal the consequences for labor. Diving deeper into how suppliers manage to maintain the required demand detail out how they employ regular workers to secure quality and consistency of production and irregular workers to cope with fluctuating orders and downward price/cost pressures. Their findings on the Moroccan garment case show even though companies might take complex taxes and economically upgrade they rely on a periphery of irregular, informal workers (often young, female migrants) for low-skilled, time-sensitive tasks (packing, loading) under poor conditions \parencite{barrientos2011economic}. Far from being a separate 'traditional' sector which was destined to phase out, informality has remained as a functional requirement and direct outcome of the competitive dynamics at the core of buyer-driven, globally fragmented production. The GVC framework thus reveals informality as constitutive of, not external to, modern capitalism."


\subsection{Debates on Informality in India}

The Indian context provides particularly rich terrain for examining formal-informal linkages due to its distinctive labor regulatory framework. India inherited a comprehensive system of protective labor legislation from the colonial and early post-independence periods, including the Industrial Disputes Act (1947), Factories Act (1948), and Contract Labour Act (1970). However, enforcement of these laws varies dramatically across states and over time, creating quasi-experimental variation in institutional environments.
Besley and Burgess (2004) exploit amendments to the Industrial Disputes Act across Indian states between 1958-1992, coding them as "pro-worker," "pro-employer," or "neutral." They find that pro-worker amendments are associated with lower output, employment, investment, and productivity in registered manufacturing, with the effects concentrated in labor-intensive industries. Importantly, they find that restrictive labor regulations increase informal sector employment, suggesting substitution between formal and informal labor.
Ahsan and Pagés (2009) extend this analysis using a continuous measure of labor regulation stringency. They confirm that stricter regulations reduce formal manufacturing growth but provide evidence of heterogeneous effects: small firms are disproportionately affected, consistent with fixed compliance costs. Bhattacharjea (2006) cautions that the Besley-Burgess coding scheme conflates amendments that affect workers' collective bargaining rights with those affecting firms' adjustment costs, potentially obscuring the mechanisms at work.
More recent work has focused on enforcement rather than statutory regulations. Chaurey (2015) shows that increases in labor inspections reduce formal employment but do not affect output or productivity, suggesting that firms respond by shifting toward contract labor or outsourcing. Adhvaryu, Chari, and Sharma (2013) find that firms located in areas with more frequent inspections are more likely to engage in "spurious compliance"—recording workers as contract laborers to evade permanent employment protections.
Critically, this literature establishes that both regulation and enforcement matter, but it largely treats informality as a labor supply response (workers pushed into self-employment) rather than examining how formal firms strategically utilize informal suppliers to circumvent regulations. My dissertation addresses this gap by modeling the vertical linkage between formal firms' outsourcing decisions and informal producers' bargaining power.


\section{Empirical Evidence on Formal-Informal Subcontracting Linkages}
\subsection{Global Evidence on Outsourcing and Labor Market Outcomes} 
A substantial empirical literature examines how firms' make-or-buy decisions affect labor market outcomes, though much of this work focuses on developed economies or international offshoring rather than domestic formal-informal linkages.
Abraham and Taylor (1996) document the rise of outsourcing in U.S. manufacturing, finding that firms substitute purchased services for direct employment in response to demand volatility and union pressure. Autor (2003) shows that temporary help employment grew as firms sought flexibility, with wages for temporary workers 20\% below comparable permanent employees. Dube and Kaplan (2010) find that outsourcing of food and janitorial services in U.S. hospitals reduces wages for affected workers by 4-7\%, with the effect mediated through weaker bargaining power rather than productivity differences.
Goldschmidt and Schmieder (2017) exploit German administrative data to show that domestic outsourcing accounts for approximately 10\% of the rise in wage inequality between 1985-2008. They find that outsourced workers earn 10-15\% less than in-house workers in identical occupations, even after controlling for match quality. Crucially, they demonstrate that wage losses persist for individual workers who transition from in-house to outsourced employment, ruling out pure selection effects.
Studies of international offshoring yield related insights. Harrison and McMillan (2011) find that manufacturing outsourcing to low-wage countries reduces domestic manufacturing wages and employment in the U.S., with effects concentrated among less-educated workers. Antràs, Fort, and Tintelnot (2017) develop a quantitative model of multinationals' sourcing decisions, showing that the elasticity of substitution between domestic and foreign inputs critically determines the distributional effects of globalization.
\subsubsection{Subcontracting in Developing Country Manufacturing}
The literature on domestic subcontracting in developing countries is more limited but growing. Basu, Chau, and Kanbur (2010) develop a theoretical model where formal firms subcontract to informal firms that employ child labor, showing how global supply chains can inadvertently perpetuate exploitative practices. Fernandez-Stark, Frederick, and Gereffi (2011) document the "putting-out" system in the apparel industry across Asia and Latin America, where formal firms supply raw materials to home-based workers who perform labor-intensive tasks for piece-rate wages well below minimum wage.
Marjit (2003) and Kar and Marjit (2009) model formal-informal linkages in the Indian context as a nested production structure where formal firms use informal services as intermediate inputs. They show theoretically that trade liberalization can increase demand for informal services if formal manufacturing expands, but they do not empirically test the mechanism or examine how bargaining power affects surplus distribution.
Recent empirical work provides more direct evidence. Nataraj (2011) uses Indian plant-level data to show that formal manufacturing firms increase outsourcing following trade liberalization, particularly in states with restrictive labor regulations. She finds that informal manufacturing employment rises in these states, but does not examine wage effects. McCaig and Pavcnik (2018) study Vietnam's enterprise formalization reforms, finding that provinces with more rapid formalization exhibit slower wage growth in household enterprises, consistent with monopsony power by newly formal firms.
Critical Gap: While this literature establishes that formal firms do outsource to informal producers and that this outsourcing is sensitive to labor regulations, it does not quantify the distribution of surplus between formal buyers and informal suppliers. Is the low profitability of informal firms explained by low productivity (as La Porta and Shleifer argue) or by extractive contract terms (as the adverse incorporation thesis suggests)? My dissertation directly addresses this question by structurally estimating the bargaining parameter $\mu$.
\subsubsection{The Indian Textile Value Chain: Institutional Context}
India's textile and garment sector provides an ideal setting for examining formal-informal linkages. The sector accounts for approximately 7\% of manufacturing output and 12\% of export earnings, with employment of over 45 million workers (Ministry of Textiles, 2023). Crucially, production is organized through extensive vertical disintegration: large formal mills and exporters subcontract labor-intensive tasks (embroidery, stitching, dyeing, finishing) to networks of small unincorporated workshops and home-based workers.
Tewari (1999) provides rich ethnographic evidence on the organization of the Ludhiana knitwear cluster in Punjab, documenting how formal exporters coordinate production across hundreds of small job-workers without formal contracts. Kadiyali, Chintagunta, and Vilcassim (2000) analyze productivity spillovers in the Tiruppur garment cluster in Tamil Nadu, finding limited technology transfer from formal to informal units despite dense network ties.
Soundararajan (2019) examines the effect of the Mahatma Gandhi National Rural Employment Guarantee Act (NREGA) on informal manufacturing wages. She finds that informal wages rise in districts with more intensive NREGA implementation, consistent with improved outside options for informal workers. However, she also finds evidence of declining informal employment, suggesting potential displacement effects. This finding underscores the importance of understanding the elasticity of substitution ($\rho$) between formal and informal inputs—a central parameter in my structural model.
Singh (2023) provides the most directly relevant recent work, examining contract labor in Indian manufacturing using ASI data. He finds that formal firms increase use of contract labor in response to pro-worker amendments to labor laws, and that contract workers earn 20-30\% less than permanent workers in observably similar jobs. However, Singh focuses on internal contract labor (workers employed on-site through contractors) rather than external outsourcing to separate enterprises. My study complements his work by analyzing Block F (job work done by others) rather than Block E (contract labor), isolating the inter-firm bargaining channel.

\section{ Labor Institutions, Enforcement, and Bargaining Power}
\subsection{Theoretical Frameworks: Institutions and Rent-Sharing}
My identification strategy rests on the theoretical premise that labor institutions affect workers' bargaining power and hence their share of economic rents. This hypothesis has deep roots in labor economics, though formalization varies across approaches.
In search-and-matching models (Mortensen and Pissarides, 1994; Pissarides, 2000), wages are determined by Nash bargaining between workers and firms over the surplus generated by the match. The worker's bargaining power parameter reflects institutional factors such as employment protection, unionization, and the generosity of unemployment insurance. Pissarides and Vallanti (2007) show that stricter employment protection increases workers' bargaining power by reducing the firm's threat point (the value of terminating the match).
Rent-sharing models in the efficiency wage tradition (Shapiro and Stiglitz, 1984; Blanchflower, Oswald, and Sanfey, 1996) provide an alternative microfoundation. If firms pay above-market wages to elicit effort or reduce turnover, then institutions that raise the cost of dismissal or strengthen workers' outside options will tilt rent-sharing toward workers. Card, Devicienti, and Maida (2014) find that Italian firms share 15-30\% of productivity gains with incumbent workers, with the pass-through higher in unionized plants.
Most relevant to my setting is the literature on bargaining in vertical supply chains. Hart and Moore (1990) and Grossman and Hart (1986) develop property rights theories where asset ownership determines parties' threat points in ex-post bargaining. Antràs (2003) applies this framework to international trade, showing that multinational firms' integration decisions depend on the hold-up problem: if suppliers make relationship-specific investments, they risk expropriation if the buyer has all the bargaining power.
Application to Informal Outsourcing: In the Indian textile context, informal workers make implicit investments (learning specific embroidery techniques, purchasing specialized equipment) that are more valuable to their current formal buyer than to alternative buyers. If labor enforcement is weak, informal workers' threat point is low—they cannot credibly walk away to formal employment at minimum wage. Thus, weak enforcement endogenizes low bargaining power ($\mu$), which I test empirically.
\subsection{Empirical Evidence on Enforcement and Wage Determination}
While the theoretical link between institutions and bargaining is well-established, empirical identification is challenging because institutional quality is endogenous to economic development. Three identification strategies have proven fruitful:
1. Cross-country comparisons with institutional instrumentation. Botero et al. (2004) code labor regulations for 85 countries and find that stricter regulations are associated with lower employment and larger informal sectors. However, these regressions suffer from omitted variable bias (institutional quality correlates with rule of law, education, etc.) and cannot isolate causal effects.
2. Within-country policy changes. Aghion, Algan, Cahuc, and Shleifer (2010) exploit French reforms to workplace representation, finding that regions with stronger unions exhibit higher pass-through of productivity to wages. Lamadon, Mogstad, and Setzler (2022) use matched employer-employee data from Sweden to estimate firm-specific rent-sharing parameters, finding substantial heterogeneity: workers in high-rent-sharing firms capture 30-40\% of productivity gains, while those in low-rent-sharing firms capture less than 10\%.
3. Spatial discontinuities in enforcement. Bhorat, Kanbur, and Stanwix (2017) exploit minimum wage variation across sectors in South Africa, finding compliance rates of 60-80\% even in sectors with weak unionization, suggesting that enforcement matters. Almeida and Carneiro (2012) show that Brazilian firms located closer to labor ministry offices face more frequent inspections and exhibit higher compliance with minimum wage laws, with effects concentrated among less-educated workers.
Application to Indian States: The Besley-Burgess (2004) framework treats state-level labor law amendments as quasi-exogenous shocks, though Bhattacharjea (2006) and others have questioned this assumption. More promisingly, variation in enforcement intensity—inspector-to-firm ratios, inspection frequencies, dispute resolution patterns—offers a plausibly exogenous source of variation, particularly in cross-sectional analysis where I control for state fixed effects. My dissertation extends this literature by testing whether enforcement affects not just wage levels but the elasticity of informal wages with respect to formal productivity—the signature of bargaining power.

\section{Methodological Approaches: Estimating Bargaining Power and Technology}
\subsection{Reduced-Form Approaches: Pass-Through Regressions}
The canonical empirical approach to identifying rent-sharing regresses worker wages on firm productivity or profitability (Van Reenen, 1996; Abowd, Kramarz, and Margolis, 1999). If labor markets were perfectly competitive, wages would depend only on workers' marginal product, which reflects individual human capital. Any correlation between firm-level productivity and individual wages (conditional on worker characteristics) indicates rent-sharing.
Card, Heining, and Kline (2013) implement this approach with German administrative data, using the Abowd-Kramarz-Margolis (AKM) decomposition to separate worker and firm wage components. They find that 20\% of wage variance is attributable to firm fixed effects, and that high-productivity firms pay higher wages even to observationally identical workers. Lamadon et al. (2022) extend this framework to allow for heterogeneous rent-sharing parameters, estimating that workers' bargaining power varies from 0.1 to 0.5 across firms.
Limitations for the Informal Sector Context: This approach requires matched employer-employee panel data to control for worker composition through individual fixed effects. Indian informal sector data (ASUSE) provides only cross-sectional establishment-level observations. I therefore adapt the pass-through framework to district-industry aggregates, testing whether the interaction between formal productivity and state enforcement predicts informal wages—a difference-in-differences style specification that allows for time-invariant state heterogeneity.
